\documentclass[]{article}
\usepackage{kk}
%opening
\title{Zasady komunikowania się w działalności naukowej i dydaktycznej}
\author{Kamil Kaczmarek}

\begin{document}

\maketitle

%automatyka, elektronika, elektrotechnika i technologie kosmiczne
%Automation, Electronics, Electrical Engineering, and Space Technologies

\section{2026-10-23}
\noindent Można robić zdj
\begin{itemize}
	\item Warto wrzucać cytat związany z tematyką danej prezentacji aby lepiej dotrzeć do widowni - zwykle pierwszy wykład
	\item Kontakt - z osobą nie, instytucją
	\item Dokładny kontakt nie tylko odnośnik - dokładne informacje 
	\item Warunki - określenie zasad współpracy (zaliczenia, przerwy, punktualność, obecność)
	\item Centrum nowoczesnej edukacji PG - na koniec miesiąca wykłady
	\item Relacje z grupą - ważne
	\item Mówić językiem budjącym a nie krytycznym
	\item Wprowdzić do tematyki zajęć
	\item Poprosić o kontakt - ograniczać telefony
	\item Poznać grupę - zapytać grupę o nich - w przypadku mniejszych grup (Analiza audytorium)
	\item Dobrowolska - jak budować silną psychikę

\end{itemize}
Pokolenie Z(1995-2014)
\begin{itemize}
	\item Brak etosu siwego włosa
	\item Potrzeba budowania autoretytetu - ważna otwartość, pasja i praktyka
	\item Komunikacja oparta na intesywnych bodźcach
	\item Bardzo ważna jest komunikacja niewerbalna
	\item Pierwsze 15min jest decydujące - buduje lub nie boduje motywacji
	\item Efekt pierwszego wrażenia:
	\begin{itemize}
		\item Efekt aureolii (hallo effect) - Galatei
		\item Efekt diabelski (rogów) - Golema
	\end{itemize}
	\item Jak wywrzeć dobre pierwsze wrażenie:
	\begin{itemize}
		\item Unikaj podnoszenia brwi podczas pierwszego spotkania
		\item Odzwierciedlanie zachowania - aby nie było sztuczne
		\item Uśmiechaj się
	\end{itemize}
	\item Dobry początek np. zdj, cytat, konrowersyjne pytanie itp. 
	\item Budowanie atmosfery:
	\begin{itemize}
		\item Aktywność uczestników
		\item Pouczicie bezpieczeństwa - możliwość popełnienia błedów
		\item Postępy procesu uczenia się - poczucie doświadczenia sukcesów
		\item Akceptacja - różne poglądy w grupie - zadbać o akceptację
	\end{itemize}
\end{itemize}

\end{document}
